% Abridged version, Section 5, part 1 of 2: the model as topological recursion, the leg dictionary, the lattice theorem.
% Continued in part 2 (L-regularity, the Gap Theorem, Theorem V, the Model Lemma).
\section{The model: topological recursion on the model curve}
\label{sec:mod}

Topological recursion is applied here, and only here, to the curve $x=\Wv(p)$, $y=p$ on $\bbP^1$, $v\neq0$, and to its rescalings; nothing about the quintic enters. This section follows \Lr{sec:mod}; the inputs (Q) are those of \cref{sec:int-inputs}.

%-------------------------------------------------------------------------------
\subsection{The model as topological recursion}
\label{sec:mod-tr}

\emph{Conventions} \Lr{conv:mod-kernel,def:mod-curve,def:mod-pair,def:mod-xi,def:mod-ancvec}. For meromorphic functions $x,y$ on $\bbP^1$ such that $\dd x$ has only simple zeros, the \emph{branch points} $\varsigma_b$, and $\dd y\neq0$ there, put $\omega_{0,1}:=y\,\dd x$ and $\omega_{0,2}(p_1,p_2):=\dd p_1\,\dd p_2/(p_1-p_2)^2$; the correlators $\omega_{h,n}$, $2h-2+n>0$, are defined by the recursion of \cite[(2-26)--(2-29)]{Dunin-Barkowski:2015caa}, a sum of residues at the branch points \Lr{conv:mod-kernel}; for $h\geq2$ the free energy is defined by $(2-2h)\FTR_h:=\sum_b\Res_{\varsigma_b}\Phi_{y\dd x}\,\omega_{h,1}$, $\dd\Phi_{y\dd x}=y\,\dd x$. On the \emph{model curve} $x=\Wv(p)$, $y=p$ one has $\dd x=(p^5-v)\,\dd p/p^2$; the branch points are $\varsigma_b:=\zeta^{-b}\varsigma$, and $p=0$, $p=\infty$ are the poles of $x$. The \emph{leg class} is $\cL:=\bbQ[v][p]+\bbQ(v)[x]$; for legs $\phi_1,\dots,\phi_n\in\cL$ the \emph{residue pairing} is
\[
\langle\phi_1\cdots\phi_n\rangle^{\mathrm{TR}}_h:=\sum_{b_1,\dots,b_n}\Res_{p_1=\varsigma_{b_1}}\cdots\Res_{p_n=\varsigma_{b_n}}\prod_{l=1}^n\phi_l(p_l)\,\omega_{h,n}(p_1,\dots,p_n),
\]
and the \emph{log leg} $-\log p$ is paired by integration along a ray $\mathfrak r$ from $0$ to $\infty$ avoiding the branch points, $\langle\cdots,-\log p,\cdots\rangle^{\mathrm{TR}}_h:=\int_{p'\in\mathfrak r}\omega_{h,n}(\dots,p',\dots)$. Let $D_x:=\dd/\dd x$, and for $a\in\bbZ/5$ and $d\geq0$ let
\[
\xi_a:=\varepsilon'_a\,v^{\alpha_a}\frac{p^{m_a}}{p^5-v},\quad(m_a,\varepsilon'_a,\alpha_a)=(4,-1,0),\ (0,1,1),\ (1,1,0),\ (2,-1,0),\ (3,-1,0),\qquad\xi_{a,d}:=D_x^d\xi_a ;
\]
for a leg $\phi$ put $q_\phi(z):=\sum_{a,d}\langle\phi,\dd\xi_{a,d}\rangle H^az^d$, where $\langle\phi,w\rangle:=\sum_b\Res_{\varsigma_b}\phi\,w$.

\begin{thm}[The model as topological recursion; {\Lr{thm:mod-tr}}]\genusstatus{AI-proposed, unchecked}
\label{thm:mod-tr}
Let $v\neq0$, $n\geq1$ and $2h-2+n>0$. The correlators of the model curve are
\[
\omega_{h,n}(p_1,\dots,p_n)=(-5)^{h-1}\sum_{a_1,\dots,a_n}\sum_{d_1,\dots,d_n\geq0}\Big\langle\prod_{i=1}^n\psi^{d_i}H^{a_i}\Big\rangle_h(v)\,\prod_{i=1}^n\dd\xi_{a_i,d_i}(p_i),
\]
a finite sum, the brackets being the ancestor correlators of $\Omega^{\mathrm{mod}}_v$. Equivalently, for legs $\phi_1,\dots,\phi_n$ in $\cL$ or equal to $-\log p$, $\langle q_{\phi_1}\cdots q_{\phi_n}\rangle_h=(-5)^{1-h}\langle\phi_1\cdots\phi_n\rangle^{\mathrm{TR}}_h$.
\end{thm}

\begin{sketch}[Sketch of the proof {\Lr{lem:mod-laplace,lem:mod-scale,prop:mod-Rtr}}]
The quoted theorem \cite[Theorem~4.1]{DuninBarkowski:2012bw}, as recalled in \cite[\S2.6]{Dunin-Barkowski:2015caa}, says: for a compact curve with simple branch points, if the matrix $R$ defined by the Laplace transforms at the branch points of the forms $\dd\xi^i$ attached to them is symplectic, satisfies \cite[(2-32)]{Dunin-Barkowski:2015caa}, and the expansion of $\omega_{0,2}$ at the branch points is the edge of $R$, then the correlators are the graph sum of $R$ on a topological field theory, paired with the forms $\dd(D_x^d\xi^i)$. The three hypotheses are verified for the curves $x=\kappa\Wv$, $y=p$: the Laplace transforms reduce to the series of the forms $p^m\dd p$ at each branch point, which satisfy (Mod1) and (Mod2), hence the T-system of \cref{prop:set-modrows}; by its uniqueness, $R_{\mathrm{TR}}(\psi)=\Mo\bigl(\psi/(\kappa\varsigma^4)\bigr)$, which is symplectic (\cref{def:set-M0}). The theorem is applied at $\kappa=4$ and transferred to $\kappa=1$ by scaling; passing to the normalised canonical basis of $\TFT^{\mathrm m}_v$ produces the factor $(-5)^{h-1}$ (the metric of topological recursion is $-\frac15$ times that of the model), and passing to the flat basis by $\Psi^{\mathrm m}_v$ turns the forms $\xi^i$ into the $\xi_a$.
\end{sketch}

%-------------------------------------------------------------------------------
\subsection{The leg dictionary and the lattice theorem}
\label{sec:mod-legs}

Let $\phi^b_0:=p,\ \log p,\ v/p,\ p^3,\ p^2$ for $b=0,\dots,4$, and $\phi^b_{k+1}:=I(\phi^b_k)$, where $I(F)$ is the termwise antiderivative of $-F\,\dd x$ without constant term; so $D_x\phi^b_{k+1}=-\phi^b_k$, and $\phi^b_k$ is homogeneous of weight $\nu_b+1+4k$ for the weights $p:1$, $v:5$, $x:4$.

\begin{prop}[The leg dictionary; {\Lr{prop:mod-LD}}, abridged]\genusstatus{AI-proposed, unchecked}
\label{prop:mod-LD}
Let $\varrho:=(-1,1,-1,-3,-2)$. For all $a,b\in\bbZ/5$ and $k,d\geq0$, $\langle\phi^b_k,\dd\xi_{a,d}\rangle=-\varrho_b\,\delta_{ab}\,\delta_{dk}$. Hence $q_{\phi^b_k}=-\varrho_b\,z^kH^b$: the leg of the insertion $v^m\psi^kH^b$ is $-\varrho_b^{-1}v^m\phi^b_k$. In particular the unit $H^0$ has the leg $y=p$; $\Phi_{y\dd x}=p^5/5-v\log p$ has $q_{\Phi_{y\dd x}}=-z\,H^0$; $-\log p$ has $q_{-\log p}=H$; the Euler--dilaton vector $\se=z\,H^0+vH$ has the leg $-p^5/5$; and $q_{f(x)}=0$ for $f\in\bbQ(v)[x]$.
\end{prop}

Functions of $x$ pair to zero with every $\omega_{h,n}$, because the sum of $\omega_{h,n}$ over the five sheets of $x$ vanishes \Lr{lem:mod-xfun}. With the dilaton equation the dictionary gives $\Fm_h=(-5)^{1-h}\FTR_h$ for $h\geq2$ \Lr{cor:mod-trF}.

The \emph{leg map} is the $\bbQ$-linear map $\Lambda(v^m\psi^kH^b):=\varrho_b^{-1}v^m\phi^b_k$ on the free $\bbQ[v]$-module on the $\psi^kH^b$, graded by the weight spaces $V_w$; it takes values in $\mathsf M:=\bbQ[p,\bar p]\oplus\log p\cdot\bbQ[v,x]$, where $\bar p:=v/p$ and $\bbQ[p,\bar p]=\bbQ[v][p]+\bbQ[v][x]$ is the ring of the Laurent polynomials whose terms $v^mp^j$ have $j\geq-m$. The element of $\bbQ[v,x]$ multiplying $\log p$ is the \emph{$\log p$-coefficient}.

\begin{prop}[Log coefficients of the legs; {\Lr{prop:mod-dii}}]\genusstatus{AI-proposed, unchecked}
\label{prop:mod-dii}
For all $b$ and $k$, the $\log p$-coefficient of $\phi^b_k$ is $\varrho_b\sum_{i=0}^k\Pi_i(\psi^kH^b)\,(-x)^i/i!$. Hence, for $X\in V_w$, the $\log p$-coefficient of $\Lambda(X)$ is $\sum_{i\geq0}\Pi_i(X)(-x)^i/i!$, and $\Lambda(X)$ is log-free if and only if $X\in\sKo$.
\end{prop}

\begin{sketch}
The $\log p$-coefficients of the $\phi^b_k$ are generated by the residues $\Res_{p=0}\Wv^n\,\dd\phi^b_0/n!$, and these are $\varrho_b\,\Pi_0(\psi^nH^b)$, since $\IC_m=\Res_{p=0}p^m\re^{-\Wv/z}\dd p$ \Lr{lem:mod-genfn,lem:mod-resid}.
\end{sketch}

Thus the vanishing tower, defined by the residue at the pole $p=0$ of the model, measures exactly the logarithms of the legs.

\begin{thm}[The lattice theorem; {\Lr{thm:mod-lattice,cor:mod-Ksplit}}, abridged]\genusstatus{AI-proposed, unchecked}
\label{thm:mod-lattice}
For every $w\geq0$:
\begin{enumerate}
\item[(a)] for $X\in\sKo\cap V_w$, $\Lambda(X)$ lies in $\bbQ[p,\bar p]=\bbQ[v][p]+\bbQ[v][x]$, exactly and not only modulo functions of $x$;
\item[(b)] $\Lambda\colon\sKo\cap V_w\to(\cL/\bbQ(v)[x])_w$ is bijective;
\item[(c)] $\sK\cap V_w=(\sKo\cap V_w)\oplus\bbQ v^{w/5}H$ if $5\mid w$, and $\sK\cap V_w=\sKo\cap V_w$ otherwise; for $X\in\sK\cap V_w$, $\Lambda(X)=\Lambda(X-\Pi_0(X)H)+\Pi_0(X)\log p$ with $\Lambda(X-\Pi_0(X)H)\in\bbQ[p,\bar p]$.
\end{enumerate}

\end{thm}

Part (a) is \cref{prop:mod-dii}; the bijectivity is proved modulo $v$ and lifted by a graded Nakayama argument \Lr{lem:mod-modv}. So an insertion from $\sK$ has a leg that is a polynomial in $p$ with coefficients in $\bbQ[v]$, plus a function of $x$, which pairs to zero, plus $\Pi_0(X)$ times the log leg.

% Abridged version, Section 5, part 2 of 2: L-regularity, the Gap Theorem, Theorem V, the Model Lemma on K.

%-------------------------------------------------------------------------------
\subsection{L-regularity, the Gap Theorem and Theorem V}
\label{sec:mod-Lreg}

The next two statements use the closed formula \cite[Theorem~7.1, (7.9), (7.10)]{ABDKS25} for the differentials of topological recursion on a rational curve $(x=x(z),y=z)$; its assumptions hold for $x=\Wv$, $y=p$, $v\neq0$.

\begin{thm}[L-regularity; {\Lr{lem:mod-abdks,thm:mod-Lreg}}]\genusstatus{AI-proposed, unchecked}
\label{thm:mod-Lreg}
For $2h-2+n>0$ there are an integer $K\geq0$ and $N_{h,n}\in\bbQ[v][p_1,\dots,p_n]$ with $\omega_{h,n}=N_{h,n}(p;v)\prod_i\dd p_i\big/\prod_i(p_i^5-v)^K$. Consequently, for legs $\phi_1,\dots,\phi_n\in\bbQ[v][p]$ the residue pairing $\langle\phi_1\cdots\phi_n\rangle^{\mathrm{TR}}_h$ lies in $\bbQ[v]$; for legs in $\cL=\bbQ[v][p]+\bbQ(v)[x]$ it is the restriction of an element of $\bbQ[v]$, and the components in $\bbQ(v)[x]$ pair to zero.
\end{thm}

\begin{sketch}
The formula of \cite{ABDKS25} writes $\omega_{h,n}$ as a finite sum of terms in $\bbQ[v][p_i,p_i^{-1},(p_i^5-v)^{-1},(p_i-p_j)^{-1}]\prod_i\dd p_i$; there are no poles at $p_i=0$ or on the diagonals, so these denominators cancel in $\bbQ[v][p]$, without negative powers of $v$. For polynomial legs the sum of the residues at the branch points is minus the residue at $p=\infty$, and $\Res_{p=\infty}p^a\,\dd p/(p^5-v)^K$ is an integer multiple of a non-negative power of $v$.
\end{sketch}

\begin{thm}[The Gap Theorem; {\Lr{prop:mod-period,thm:mod-gap}}]\genusstatus{AI-proposed, unchecked}
\label{thm:mod-gap}
On the model curve, for every $h\geq1$, $\omega_{h,1}=\dd\Phi_h$ with $\Phi_h$ rational and holomorphic at $p=0$ and $p=\infty$, and $\int_\infty^0\omega_{h,1}=\frac{B_{2h}}{2h}\,v^{1-2h}$ along any path avoiding the branch points. For every $h\geq2$,
\[
\FTR_h=\frac{B_{2h}}{2h(2h-2)}\,v^{2-2h},\qquad\text{hence}\qquad\Fm_h=c_hv^{2-2h},\quad c_h=(-5)^{1-h}\frac{B_{2h}}{2h(2h-2)}=\frac{(-1)^{h-1}B_{2h}}{2h(2h-2)\,5^{h-1}} .
\]
\end{thm}

\begin{sketch}
By \cite[(7.10)]{ABDKS25}, $\omega_{h,1}$ is the differential of an explicit rational function $\Phi_h$, a finite sum of terms $D_x^{r-1}(a_r/x')$. It vanishes at $p=\infty$ by a count of degrees. At $p=0$ only the leading parts of the factors contribute, those of the pole $v/p$ of $\Wv$, and the value is $v^{1-2h}$ times a universal constant, which is $B_{2h}/(2h)$. The free energy follows from $(2-2h)\FTR_h=\langle\Phi_{y\dd x}\rangle^{\mathrm{TR}}_h$ with $\Phi_{y\dd x}=p^5/5-v\log p$, whose polynomial part pairs to zero.
\end{sketch}

For logarithmic legs the tool is the variational formula of topological recursion \cite[Theorem~5.1, (5-14)]{Eynard:2007kz}: for points $J$ off the branch points, $n\geq1$ and $h+n>1$,
\[
\partial_v|_x\,\omega_{h,n}(J)=\int_{q\in\mathfrak r}\omega_{h,n+1}(J,q),
\]
where $\partial_v|_x$ is the derivative in $v$ at fixed $x$ \Lr{thm:mod-var}. With $\partial_v|_x=\partial_v|_p-\mathrm{Lie}_{\cX}$, $\cX:=\frac p{p^5-v}\frac{\dd}{\dd p}$, and an integration by parts, it becomes a recursion on the number of log legs.

\begin{thm}[Theorem V; {\Lr{thm:mod-V}}]\genusstatus{AI-proposed, unchecked}
\label{thm:mod-V}
Let $h\geq0$, $a,k\geq0$, $n=a+k\geq1$, $2h-2+n>0$, and $\phi_1,\dots,\phi_a\in\bbQ[v][p]$.
\begin{enumerate}
\item If $a\geq1$, then $\langle\phi_1\cdots\phi_a,(-\log p)^k\rangle^{\mathrm{TR}}_h\in\bbQ[v]$.
\item If $a=0$, then $\langle(-\log p)^k\rangle^{\mathrm{TR}}_h=(\FTR_h)^{(k)}(v)$, with $\FTR_h$ of \cref{thm:mod-gap} for $h\geq2$, $\FTR_1:=-\frac1{12}\log v$ (here $k\geq1$) and $\FTR_0:=\frac12v^2\log v$ (here $k\geq3$).
\item The recursion (V1) holds:
\[
\begin{aligned}
\langle\phi_1\cdots\phi_a,(-\log p)^{k+1}\rangle^{\mathrm{TR}}_h={}&\partial_v\langle\phi_1\cdots\phi_a,(-\log p)^k\rangle^{\mathrm{TR}}_h-\sum_{i=1}^a\langle\cdots,\partial_v\phi_i,\cdots,(-\log p)^k\rangle^{\mathrm{TR}}_h\\
&+\sum_{i=1}^a\langle\cdots,(\cX\dd\phi_i)_+,\cdots,(-\log p)^k\rangle^{\mathrm{TR}}_h ,
\end{aligned}
\]
where $(\cX\dd\phi)_+\in\bbQ[v][p]$ is the quotient in the division of $p\,\phi'(p)$ by $p^5-v$.
\end{enumerate}
\end{thm}

\begin{sketch}
(iii) is the variational formula, differentiated under the pairings. For (i), induct on $k$: the right side of (V1) consists of a $v$-derivative of, and pairings with, $a$ polynomial legs and $k$ log legs, and the induction starts with \cref{thm:mod-Lreg} or, in genus $0$, with the explicit $\omega_{0,3}$. For (ii), (V1) with $a=0$ reads $\langle(-\log p)^{k+1}\rangle^{\mathrm{TR}}_h=\partial_v\langle(-\log p)^k\rangle^{\mathrm{TR}}_h$, so it suffices to treat the first admissible $k$: for $h\geq2$ this is \cref{thm:mod-gap}, and for $h=0,1$ it follows from $\omega_{0,3}$ and from a local formula for $\omega_{1,1}$.
\end{sketch}

%-------------------------------------------------------------------------------
\subsection{The Model Lemma on \texorpdfstring{$\sK$}{K}}
\label{sec:mod-MLK}

\begin{thm}[The Model Lemma on $\sK$; {\Lr{thm:mod-MLK}}]\genusstatus{AI-proposed, unchecked}
\label{thm:mod-MLK}
$\ML{h}$ holds for every $h\geq0$.
\end{thm}

\begin{proof}[Proof, given the results of this section]
Part (ii) of $\ML{h}$ is \cref{thm:mod-gap}. For part (i), let $X_i\in\sK\cap V_{w_i}$, $i=1,\dots,n$. By \cref{thm:mod-lattice}~(c), $X_i=Y_i+\Pi_0(X_i)H$ with $Y_i\in\sKo\cap V_{w_i}$ and $\Pi_0(X_i)\in\bbQ v^{w_i/5}$, a monomial with a non-negative exponent. By \cref{prop:mod-LD}, $X_i=q_{\ell_i}$ for the leg $\ell_i:=-\Lambda(Y_i)+\Pi_0(X_i)\,(-\log p)$, and by \cref{thm:mod-lattice}~(a), $-\Lambda(Y_i)=\phi_i+f_i(x)$ with $\phi_i\in\bbQ[v][p]$ and $f_i\in\bbQ[v][x]$. By \cref{thm:mod-tr}, multilinearity and the vanishing of the pairings with functions of $x$,
\[
\langle X_1\cdots X_n\rangle_h=(-5)^{1-h}\sum_{S\subseteq\{1,\dots,n\}}\prod_{i\notin S}\Pi_0(X_i)\,\big\langle(\phi_i)_{i\in S},(-\log p)^{n-|S|}\big\rangle^{\mathrm{TR}}_h .
\]
Every term with $S\neq\emptyset$ lies in $\bbQ[v]$, by \cref{thm:mod-Lreg} if $S=\{1,\dots,n\}$ and by \cref{thm:mod-V}~(i) otherwise. By \cref{thm:mod-V}~(ii) and $\Fm_h=(-5)^{1-h}\FTR_h$ (for every $h\geq0$, with $\FTR_0$, $\FTR_1$ as in \cref{thm:mod-V}), the term $S=\emptyset$ is $(\Fm_h)^{(n)}(v)\prod_i\Pi_0(X_i)$.
\end{proof}



